\documentclass[12pt]{article}
\usepackage[utf8]{inputenc}
\usepackage{amsmath, amssymb, amsfonts}
\usepackage{geometry}
\geometry{a4paper, margin=1in}
\title{MA102 Answer Script Solution}
\author{Roll Number: 240104105}
\date{\today}
\begin{document}
\maketitle
\section*{Solution}

Q-4) Space for answering Question 4 only.

a) fog and gof are identity on V [that is they give some output as input]
[fog] = [gof] = I (Identity matrix)

This proves that f and g are inverse functions of each other.

Inverse operator of a linear operator is also linear,
(f(x+y)=f(x)+f(y)) => f^-1(x+y)=f^-1(x)+f^-1(y)
g is also a linear operator

b) Given: B = [b_1, b_2, ..., b_n]^T is a basis for V.
c_1b_1 + c_2b_2 + ... + cnbn = O [where e_1, e_2, ..., en^T R
and c_1 = ... = c_n = 0]

Applying f on both sides:
f(c_1b_1 + c_2b_2 + ... + cnbn) = f(0)
By property of linear operator [f(e_1x+e_2y+...+en^T)] = f(e_1x+f(e_2y+...+en^T)]
This implies that f(b_1), f(b_2), ..., f(b_n) are linearly independent vectors.

Now let any w E V, the function is surjective span(B) = V
c_1w + c_2b_2 + ... + cnbn = w [c_1, c_2, ..., (n+1)ER]
Applying f on both sides:
c_1fo(b_1) + c_2fo(b_2) + ... + cnfn(b_n) = fo(w)

\end{document}